\subsection{Beginnings}

\content{
        \begin{definition} Each individual which is casting a vote is called a
        \textbf{player}, which we will denote $P_1, P_2, \dots, P_N$, where $N$
        is the number of players.

        Each player is given a \textbf{weight}, which usually 
        represents how many votes they get. 

        The \textbf{quota} is the minimum weight needed for the proposal to be
        approved. 

        We will often represent all of this information in the following form: 
        $$ [q: w_1, w_2, \dots, w_N] $$
        where $q$ represents the quota, and $w_i$ is the weight of player $P_i$.
        \end{definition}
}{% presentation
\vspace{2in}
}{% nopresentation
\vspace{2in}
}{% comments
}

\content{
        \begin{example} Suppose that in a small company, There are four
        shareholders, the first has a 30\% ownership, the second and third both
        have a 25\% ownership, and the fourth a 20\% ownership. The company
        by-laws state that more than 50\% of the ownership must approve any
        decision to open a new location. Represent this situation as a wieghted
        voting system. 
        \end{example}
}{% presentation
\vspace{2in}
}{% nopresentation
\vspace{2in}
}{% comments
$$ [51: 30, 25, 25, 20] $$
}

\content{
        \begin{example} What is wrong with the following weighted voting system? 
        $$ [11: 4, 3, 2, 1] $$
        \end{example}
}{% presentation
\vspace{2in}
}{% nopresentation
\vspace{2in}
}{% comments
This is gridlock, nothing can be approved. 
}

\content{
        \begin{example} What is wrong with the following weighted voting system? 
        $$ [3: 3,2,1] $$
        \end{example}
}{% presentation
\vspace{2in}
}{% nopresentation
\vspace{2in}
}{% comments
This is anarchy, since both yes and no could reach the quota simultaneously. 
}

\subsubsection{Power}

\content{
        \begin{example} Consider the weighted voting system 
        $$ [8: 9,3,2,1]. $$
        Where does the power lie? Why?
        \end{example}
}{% presentation
\vspace{2in}
}{% nopresentation
\vspace{2in}
}{% comments
Player 1 has all the power. 
}


\content{
        \begin{definition} A player is said to have \textbf{veto power} if they
        are necessary for a proposal to pass. 
        \end{definition}

        \begin{definition} A \textbf{dictator} is a player who can reach the
        quota by themselves and also has veto power. 
        \end{definition}

        \begin{example} Determine all the players who have veto power in the
        weighted voting system
        $$ [22: 8, 6, 4, 4, 3, 2, 1]. $$
        \end{example}
}{% presentation
\vspace{2in}
}{% nopresentation
\vspace{2in}
}{% comments
In the previous example, player 1 is a dictator. \\
How do we determine who has veto power? Derive the formula 
$$ V - q < w_i. $$
}

\content{
        \begin{example} Determine all the players who have veto power in the
        following weighted voting system: 
        $$ [34: 12, 11, 10, 7, 3 ]. $$
        \end{example}
}{% presentation
\vspace{3in}
}{% nopresentation
\vspace{3in}
}{% comments
Players 1, 2, and 3 all have veto power. 
}

\content{
        \begin{definition} A player is called a \textbf{dummy} if their vote is
        never essential for a group to reach quota. 
        \end{definition}

        \begin{example} Identify any dummies in the weighted voting system 
        $$ [10: 7, 6, 2]. $$
        \end{example}
}{% presentation
\vspace{2in}
}{% nopresentation
\vspace{2in}\newpage
}{% comments
}

\content{
        Before we can define power, we need a few more definitions. 

        \begin{definition} A \textbf{coalition} is a group of players voting the
        same way. 
        \end{definition}

        \begin{definition} A coalition is said to be a \textbf{winning
        coalition} if the players in the coalition reach the quota. 
        \end{definition}

        \begin{definition} A player in a coalition is said to be critical if
        them leaving the coalition would change it from a winning coalition to a
        losing coalition. 
        \end{definition}

        \begin{example} Consider the weighted voting system: 
        $$ [14: 7, 6, 3, 3]. $$
        Identify all winning coalitions. 
        \end{example}
}{% presentation
\vspace{1in}
}{% nopresentation
\vspace{2in}
}{% comments
}

